\section[Overview]{Overview}
This exercise is a step-by-step modeling project, 
which we will progress over the first 4 days of the course at the end of each day.
Each step corresponds and is focusing on the specific process that are being
newly taught earlier on that day.      
Generally there will be detailed step-by-step instructions provided. 
As part of this project you will also learn how to compile a project-specific 
reaction database. These skills will prepare you for your own 
reactive transport modeling project in the future. The project involves 
a two-dimensional simulation problem that eventually includes

\begin{itemize} 
\item advective-dispersive transport 
\item mineral reactions 
\item cation exchange and surface complexation, and 
\item kinetically controlled redox reactions. 
\end{itemize} 

The exercise is a conceptual modeling problem, meaning that we will not compare the simulation 
results to any measured data. The purpose of the exercise is mainly to illustrate 
how different processes individually and jointly affect groundwater quality evolution.

The problem is set up as a simple vertical transect representing a small catchment with 
groundwater recharge being the only driver for groundwater flow, i.e., there is no pumping or injection.
The model extends all the way to the catchment boundary and all outflow occurs through groundwater discharge
towards a river. We will study the evolution of the groundwater quality within the aquifer and how this affects   
the composition of the water that discharges to the river. 
The simulations will mimic, although in a crude way, how groundwater quality changes may occur in response to the increased atmospheric sulfur deposition \citep{ISI:000179120500001} that has occurred in the past, and in some places still occurs. 
Those elevated atmospheric sulfur deposition  occurred in conjunction 
with air pollution from industrial emissions \citep{ISI:000266220900013,ISI:000084664000003}. 
The increased sulfur input into the groundwater system is associated with a decline of the pH within the 
rainwater (i.e., acid rain). The increased acidity of the recharge water can cause the enhanced dissolution of selected minerals \citep{ISI:000455694400029}, which act as pH buffers. The dissolution of these minerals will raise the pH of the groundwater until they are depleted. With the exercise you will compare different scenarios and explore the influence of the initial concentration of the buffering mineral on the longevity of the buffering process. Moreover, the difference in composition between native groundwater and rainwater and the associated mineral dissolution will cause release and re-sequestration of ions that are originally associated with sediment cation exchange and/or surface complexation sites. As part of the the exercise you will incorporate the relevant cation exchange and surface complexation reactions and explore their influences on the groundwater quality evolution. Lastly, the rainwater also contains dissolved oxygen, which can trigger redox reactions if the receiving aquifer contains minerals that can act as reductants. For example, if pyrite is prevalent in a reducing aquifers, the intrusion of aerobic recharge water would oxidize the pyrite and potentially cause metal(loid) mobilization. As part of this exercise you will incorporate the kinetically controlled pyrite oxidative dissolution reaction and explore the influence of pyrite oxidation on groundwater quality.
